\noindent\textbf{ZUR HANNOVERSCHEN TRIANGULATION.}\hfill347
\begin{center}\textsc{Gauss} an \textsc{Schumacher}.\quad Göttingen, 10. September 1818.\end{center}

Ich beeile mich, Ihnen mitzuteilen, dass ich von unserem Minister \textsc{Arnswaldt} den
Auftrag erhalten habe, die zur Verbindung einer hannoverschen Triangulation mit
der Ihrigen nötigen Messungen in Lüneburg vorzunehmen und dazu das
Nötige mit Ihnen zu vereinbaren. Er macht mir zugleich Hoffnung, dass
demnächst auch die Fortsetzung selbst wohl zustande kommen wird, und
es freut mich, dass diese nun durch die in Lüneburg vorzunehmenden Arbeiten
gesichert werden kann.\ .\ .\ .\ .\ .\ .

\begin{center}\textsc{Gauss} an \textsc{Schumacher}.\quad Göttingen, 20. Mai 1820.\end{center}
\begingroup
\makeatletter
\renewcommand{\@makefnmark}{\hbox{\normalfont[\textsuperscript{*})]}}
\long\def\@makefntext#1{\parindent1.5em\noindent[\textsuperscript{*})\ #1]}
\makeatother
\par\noindent\hspace*{\parindent}.\ .\ .\ .\ .\ .\ In dieser Ungewissheit\footnote{Über \textsc{Schumachers} Aufenthalt.} adressiere ich diesen Brief nach Copenhagen
und wünsche sehnlich, dass er Sie erreicht und \textls[120]{bald} erreicht. Er
soll Ihnen nämlich die Nachricht überbringen, dass infolge eines Schreibens
des Grafen \textsc{von Münster} aus London als Antwort auf meinen vor einem Jahr
\makeatletter
\renewcommand{\@makefnmark}{\hbox{\normalfont[\textsuperscript{**})]}}
\long\def\@makefntext#1{\parindent1.5em\noindent[\textsuperscript{**})\ #1]}
\makeatother
von Ihnen freundlicherweise besorgten Brief\footnote{Abgedruckt in Band IV, S.~482/483.},
\begin{quote}
»der König die Fortsetzung der Gradmessung durch das Königreich
Hannover genehmigt hat«.
\end{quote}
\par\medskip
.\ .\ .\ .\ .\ .
\endgroup

\begin{center}\textsc{Gauss} an \textsc{Schumacher}.\quad Göttingen, 4. März 1821.\end{center}

.\ .\ .\ .\ .\ .\ Ursprünglich wollte ich in Hamburg beginnen und so
von Norden nach Süden messen. Leider halten mich die
Schreibfaulheit und Unzuverlässigkeit aller Instrumentenbauer, mit denen ich zu tun
habe, so sehr auf, und ich habe bis jetzt noch gar nichts von den nötigen
Hilfsmitteln in Händen, dass die daraus entstehenden Schwierigkeiten
noch viel größer wären, wenn ich jenen Plan befolgte. Dieser und noch
\textls[120]{verschiedene andere wichtige} Gründe zwingen und bewegen mich
zum umgekehrten Plan, von Süden nach Norden zu messen.\ .\ .\ .\ .\ .\ .
\par
\makeatletter
\insert\footins{\reset@font\footnotesize\hsize\columnwidth\parindent0pt
\smallskip\hbox to\hsize{\hfil44\textsuperscript{*}}\par}
\makeatother

\clearpage
\noindent348\hfill\textbf{BRIEFWECHSEL MIT SCHUMACHER.}
\begin{center}\textsc{Gauss} an \textsc{Schumacher}.\quad Wulfsode, 18. September 1822.\end{center}

.\ .\ .\ .\ .\ .\ Soweit eine vorläufige Besichtigung des Geländes ein Urteil erlaubt,
wäre es nicht unmöglich, \textls[120]{die ganze Linie} von Breithorn bis Eschede
[Scharnhorst] (11220 Meter lang) unmittelbar zu messen. Welch eine herrliche
Basis wäre das!\ .\ .\ .\ .\ .\ .

\begin{center}\textsc{Gauss} an \textsc{Schumacher}.\quad [Göttingen, Januar 1825.]\end{center}

.\ .\ .\ .\ .\ .\ Wenn ich alle größeren und kleineren Schneisen aus den Jahren
1821---1824 zusammenzähle, von solchen, für die vielleicht ein Dutzend Bäume
gefällt wurden, bis zu den größten, so mögen es etwa 16 oder 17 Schneisen gewesen
sein. Die allergrößte, der Ausdehnung nach, lag im Becklinger
Holz unweit der Straße von Bergen nach Soltau.\ .\ .\ .\ .\ .\ .

\begin{center}\textsc{Gauss} an \textsc{Schumacher}.\quad Dangast, 1 Stunde von Varel, 20. Juni 1825.\end{center}

.\ .\ .\ .\ .\ .\ Ich selbst zeichne [beim Heliotrop] gewöhnlich durch die drei Fußspitzen
einen Kreis, dessen Mittelpunkt ich als Zielpunkt betrachte.
Bei meinen beiden neuesten Heliotropen ist auch der Mittelpunkt selbst durch eine
Spitze markiert, was sehr praktisch ist.\ .\ .\ .\ .\ .\ .

\begin{center}\textsc{Gauss} an \textsc{Schumacher}.\quad Göttingen, 14. Januar 1829.\end{center}

Da wir noch wenig darüber wissen, ob die Gestalt der Erde von der Gestalt
des mittleren Ellipsoids in längeren oder in kürzeren Undulationen abweicht,
bin ich jedenfalls der Meinung, dass man am besten immer das mittlere
Ellipsoid zugrunde legt.

Wenn ich es jedoch nicht rundweg verwerfen will, gelegentlich ein
oskulierendes Ellipsoid zugrunde zu legen, so kann ich dies doch nur dort
für zulässig halten, wo man die Mittel hat, ein oskulierendes Ellipsoid zu bestimmen,
d.\ h.\ ein solches, dessen Krümmung \textls[120]{sowohl} in Richtung des Meridians
als auch in der dazu senkrechten Richtung der Krümmung der wirklichen Gestalt so nahe wie
möglich kommt. Dies lässt sich aber nur durch die Verbindung von \textls[120]{Messungen}
in beiden Richtungen erreichen (Breiten- und Längengradmessung).\ .\ .\ .\ .\ .\ .

\clearpage
\noindent\textbf{ZUR HANNOVERSCHEN TRIANGULATION.}\hfill349
\par\medskip
\begingroup
\makeatletter
\renewcommand{\@makefnmark}{\hbox{\normalfont\textsuperscript{*})}}
\makeatother
Wenn ein Meridian nicht wirklich elliptisch ist, kann man allerdings
eine Ellipse berechnen, die sich an zwei Teilstücke eines Bogens anschließt, und
man mag diese meinetwegen eine oskulierende Ellipse nennen. Doch durch
Drehung dieser Ellipse um \textls[120]{ihre} kleine Achse\footnote{\hspace{.25em}Von ihr lässt sich nur sagen, dass sie parallel zur Erdachse ist, ohne mit ihr zusammenzufallen. In der Regel liegt sie sogar weit von ihr entfernt.} entsteht keine Fläche, die
man als oskulierendes Ellipsoid bezeichnen darf. Mit anderen Worten besteht zwischen
dem wirklichen Wert des Längengrades und demjenigen Längengrad, den man
auf dem durch Drehung jener oskulierenden Ellipse [entstandenen Ellipsoid]
berechnet, überhaupt kein Zusammenhang. Verwechselt man beide, so setzt
man sich den größten Fehlern aus.\ .\ .\ .\ .\ .\ .
\endgroup
\begin{center}\rule{.26\linewidth}{.4pt}\end{center}
